\chapter{\texorpdfstring{A Equação Unificada: Derivazione di
\(T_s^{(2)}\)}{A Equação Unificada: Derivazione di T\_s\^{}\{(2)\}}}\label{lequação-unificata-derivazione-di-t_s2}

\begin{quote}
\textbf{Argumento central:} la stabilità di un sistema egemonico è
mensurável come rapporto strutturato tra forze di coerenza e forze di
decoerência, modulato da due variáveis --- antifragilidade (\(A\)) e
involução (\(\Omega\)) --- che agiscono su dimensioni diverse del
sistema.
\end{quote}



\section{Premessa di Coerenza
Notazionale}\label{premessa-di-coerenza-notazionale}

Prima di derivare l'equação, fissiamo in modo univoco la notazione ---
che nel corpus DEQ precedente oscillava tra due convenzioni.

\Hypoth{} \textbf{Convenzione unificata adottata in questo volume:}

Il \textbf{fator de desconto} \(\delta \in [0, 1]\) è un \emph{orizzonte
temporale}: misura quanto il sistema valuta il futuro com relação al
presente. È \textbf{stabilizzante}. Nel Teorema Folk della
\emph{Geometria dell'Egemonia}, la cooperazione sostenibile richiede
\(\delta \geq \delta^* = 0{,}50\). Alto \(\delta\) = bassa preferenza
temporale = cooperazione possibile = stabilità.

\Proved{} \textbf{Conseguenza formale:} \(\delta\) appare \emph{al
numeratore} del \(T_s\). Un sistema con \(\delta\) alto è più coerente.

L'incoerenza del \emph{Memoriale Tecnico DEQ} (2026-04-23), dove
\(\delta\) compariva al denominatore, è qui risolta: quella era
un'abbreviazione tipografica della quantità complementare
\((1-\delta)\), che misura la \textbf{preferenza per il presente}
(destabilizzante). Per chiarezza assoluta, in questo volume usiamo
sempre \(\delta\) come orizzonte stabilizzante.



\section{La Doppia Eredità: Geometria e
Topografia}\label{la-doppia-eredituxe0-geometria-e-topografia}

\subsection{Il Payoff Macro della Geometria
dell'Egemonia}\label{il-payoff-macro-della-geometria-dellegemonia}

Nel volume \emph{La Geometria dell'Egemonia} (Marcelino, 2026), l'attore
\(i\) ottiene payoff di periodo:

\[
\Pi_i(t+1) = H_i \cdot Q_i \cdot (1 - \delta_i \cdot s_{i,3})
\]

onde \(H_i\) è l'impronta sistemica,
\(Q_i = (s_{i,1}+s_{i,2}+s_{i,3})/\sqrt{3}\) la qualità strategica, e
\(s_{i,3}\) la risorsa relazionale (legitimidade).

\subsection{La Funzione Predittiva Micro della Topografia del
Discorso}\label{la-funzione-predittiva-micro-della-topografia-del-discorso}

Nel volume \emph{La Topografia del Discorso Emancipativo} (Marcelino,
2026, \S{}5.7), per un singolo acoplamento emittente-ricevente:

\[
\Pi(t+1) = M_e \cdot Q_0 \cdot (1 - \delta_R \cdot z_R)
\]

onde \(M_e\) è la massa comunicativa dell'emittente,
\(Q_0 = (x+y+z)/\sqrt{3}\) la qualidade ética del discorso, \(\delta_R\)
il fator de desconto del ricevente, \(z_R\) la sua complessità
dialettica.

\subsection{\texorpdfstring{\Proved{} L'Identità Formale (Risultato
Strutturale)}{ L'Identità Formale (Risultato Strutturale)}}\label{lidentituxe0-formale-risultato-estrutural}

Le due formule \textbf{non sono analoghe --- sono la stessa equação
applicata a due scale diverse}:

{\def\LTcaptype{none} % do not increment counter
{\small\begin{longtable}[]{@{}
  >{\raggedright\arraybackslash}p{(\linewidth - 4\tabcolsep) * \real{0.3333}}
  >{\raggedright\arraybackslash}p{(\linewidth - 4\tabcolsep) * \real{0.3333}}
  >{\raggedright\arraybackslash}p{(\linewidth - 4\tabcolsep) * \real{0.3333}}@{}}
\toprule\noalign{}
\begin{minipage}[b]{\linewidth}\raggedright
Grandezza
\end{minipage} & \begin{minipage}[b]{\linewidth}\raggedright
Topografia (micro-comunicativo)
\end{minipage} & \begin{minipage}[b]{\linewidth}\raggedright
Geometria (macro-egemonico)
\end{minipage} \\
\midrule\noalign{}
\endhead
\bottomrule\noalign{}
\endlastfoot
Massa di proiezione & \(M_e\) & \(H_i\) \\
Qualità sull'asse etico & \(Q_0 = (x+y+z)/\sqrt{3}\) &
\(Q_i = (s_{i,1}+s_{i,2}+s_{i,3})/\sqrt{3}\) \\
Fattore di erosione & \((1 - \delta_R z_R)\) &
\((1 - \delta_i s_{i,3})\) \\
Soglia di cooperazione (Teorema Folk) & \(\delta^*_E = 0{,}50\) &
\(\delta^* = 0{,}50\) \\
\end{longtable}}
}

\Proved{} \textbf{Conseguenza:} la coincidenza numerica
\(\delta^* = 0{,}50\) tra Teorema Folk Discursivo (Topografia Ap.~A) e
Teorema Folk da Geometria \emph{non è una coincidenza} --- è
l'invariância de escala di un'unica struttura matematica. Questo è il
primo risultato estrutural della DEQ$^2$: la stessa legge governa il
pagamento dell'emittente verso il ricevente e il pagamento dell'egemone
verso l'attore subordinato.

\Hypoth{} \textbf{Mossa di passaggio al livello sistemico:} non
sostituiamo la logica dell'attore singolo --- la \emph{integriamo}. Il
payoff individuale \(\Pi_i(t+1)\) resta valido come legge local; ciò
che cerchiamo è la quantità che descrive lo \emph{stato collettivo},
\(T_s\), come ordine emergente dal campo \(\{\Pi_j\}\).



\section{3.1bis Il Bridge Micro/Macro via
Kuramoto}\label{bis-il-bridge-micromacro-via-kuramoto}

\Hypoth{} Considereremo il sistema come una popolazione di \(N\)
accoppiamenti elementari (cittadino-discorso, attore-egemone,
élite-massa) ciascuno con payoff instantâneo \(\Pi_j(t)\) e fase
\(\theta_j(t)\) (la posizione nel ciclo di legittimazione/decoerência).
Definiamo il \textbf{parâmetro de ordem di Kuramoto} (1984):

\[
\Phi(t) e^{i\psi(t)} = \frac{1}{N}\sum_{j=1}^{N} e^{i\theta_j(t)} \qquad \Phi \in [0, 1]
\]

\begin{itemize}
\tightlist
\item
  \(\Phi \to 1\): tutte le fasi allineate (coerenza piena);
\item
  \(\Phi \to 0\): fasi distribuite uniformemente (decoerência piena).
\end{itemize}

\Hypoth{} \textbf{Operatore di trasposizione:} definiamo la
\emph{quantità sistemica} derivata da una grandeza local \(X_j\) come

\[
\langle X \rangle_\Phi := \Phi \cdot \frac{1}{N}\sum_j X_j
\]

L'aggregazione ``naive'' \(\frac{1}{N}\sum X_j\) è recuperata per
\(\Phi = 1\); in regime decoerente \(\Phi < 1\) il sistema \emph{perde
sostanza} anche a parità di payoff individuali. È questa la differenza
tra un USA decoerente con alti payoff individuali e un Brasil
\(\Phi\)-engine con payoff individuali più modesti ma allineati.

\Proved{} \textbf{Risultato:} \(T_s^{(2)}\) derivato nelle sezioni
successive è formalmente \(\langle T_s^{\text{loc}}\rangle_\Phi\) del
campo Topografia, modulato da \(A\) e \(\Omega\). La forma sistemica
\(T_s^{(2),\text{sys}}\) del Cap.~4 esplicita questo passaggio.



\section{\texorpdfstring{O Fundamento Evolutivo de \(\Phi\): Coopera\c{c}\~{a}o Capturada}{O Fundamento Evolutivo de Phi: Cooperacao Capturada}}\label{il-fondamento-evolutivo-di-phi}

\Hypoth{} \textbf{O problema causal.} A se\c{c}\~{a}o~3.1bis introduziu
\(\Phi\) como par\^{a}metro de ordem de Kuramoto --- uma quantidade
matem\'{a}tica que mede o alinhamento de fase. A quest\~{a}o causal
permanece aberta: \emph{por que} as fases se alinham ou decoer\^{e}m?
Qual \'{e} o mecanismo que produz \(\Phi > 0\) como estado est\'{a}vel,
e a decoer\^{e}ncia como estado produzido?

A resposta vem da biologia evolutiva: a coopera\c{c}\~{a}o \'{e} o
estado natural quando os mecanismos que a sustentam superam os que a
corroem. \(\Phi \to 0\) n\~{a}o \'{e} o equil\'{i}brio de base ---
\'{e} um estado \emph{engenheirado} pela extra\c{c}\~{a}o sistem\'{a}tica.

\Proved{} \textbf{A condi\c{c}\~{a}o de Hamilton (1964).} A
coopera\c{c}\~{a}o \'{e} evolutivamente est\'{a}vel numa popula\c{c}\~{a}o
se e somente se:

\[
rB > C \tag{H}
\]

onde \(r\) \'{e} o coeficiente de parentesco, \(B\) \'{e} o benef\'{i}cio
para o receptor e \(C\) o custo para o doador. Esta condi\c{c}\~{a}o ---
derivada para a biologia e aplicada aqui como isomorfismo estrutural
(\Hypoth{}) --- tem o seu an\'{a}logo geopo\-l\'{i}tico: a
coopera\c{c}\~{a}o sist\^{e}mica \'{e} est\'{a}vel quando o benef\'{i}cio
coletivo supera o custo individual de cooperar, ponderado pela coer\^{e}ncia
estrutural do sistema.

\Hypoth{} \textbf{Correspond\^{e}ncia estrutural com \(\Phi\).}
O par\^{a}metro de ordem \(\Phi > 0\) \'{e} a assinatura sist\^{e}mica do
regime \(rB > C\): as fases se alinham quando os mecanismos cooperativos
(reciprocidade, norma institucional, confian\c{c}a) superam os mecanismos
de defec\c{c}\~{a}o. N\~{a}o \'{e} uma identidade --- \'{e} uma
correspond\^{e}ncia de regime:

\[
rB > C \;\xleftrightarrow{\;\text{\Hypoth{}}\;}\; \Phi > 0
\]

\Hypoth{} \textbf{O CPGG como mecanismo de captura.} O
\textit{Captured Public Goods Game} (CPGG --- §D.4) formaliza o mecanismo
pelo qual a coopera\c{c}\~{a}o \'{e} erodida por uma sub-rede extrativa de
fra\c{c}\~{a}o \(\phi \ll 1\) com retroalimenta\c{c}\~{a}o positiva
end\'{o}gena (la\c{c}o de captura, Proposi\c{c}\~{a}o~1 \Proved{}):

\[
|\mathcal{E}| \uparrow \;\Rightarrow\; \sigma_{\text{san}} \downarrow \;\Rightarrow\; \pi_E \uparrow \;\Rightarrow\; |\mathcal{E}| \uparrow
\]

O grau efetivo dos cooperadores em redes livre de escala (a topologia
das redes reais) colapsa com \(\phi\):

\[
k_{\text{eff}}(\phi) = \bar{k} \cdot \frac{1 - \phi^{(\gamma-2)/(\gamma-1)}}{1-\phi} \quad (\text{\Proved{} F3})
\]

O limiar cr\'{i}tico \(\phi^*\) \'{e} o ponto em que a condi\c{c}\~{a}o
\(b/c > k_{\text{eff}}(\phi)\) deixa de valer e \(\Phi\) colapsa.
O sistema sinaliza a transi\c{c}\~{a}o \emph{antes} do colapso
cooperativo: o monitoramento de \(\phi\) fornece um indicador precoce
de \(T_s^{(2),\text{sys}} \to 0\).

\Proved{} \textbf{O fundamento da sele\c{c}\~{a}o multinível.} A
justifica\c{c}\~{a}o causal da mudan\c{c}a de n\'{i}vel micro/macro
\'{e} oferecida pela \textit{teoria da sele\c{c}\~{a}o multinível}
(Wilson \& Sober 1994; Okasha 2006). A equa\c{c}\~{a}o de Price (1970)
--- uma identidade alg\'{e}brica exata --- decompõe a varia\c{c}\~{a}o
do comportamento cooperativo num termo inter-grupo (coer\^{e}ncia) e
num termo intra-grupo (extra\c{c}\~{a}o):

\[
\Delta\bar{z} = \underbrace{\mathrm{Cov}(W_k,\bar{z}_k)}_{\text{inter-grupo}} + \underbrace{\mathbb{E}_k\!\left[\mathrm{Cov}(W_{ik},z_{ik})\right]}_{\text{intra-grupo}} \tag{Price 1970}
\]

\Hypoth{} A correspond\^{e}ncia estrutural \'{e}: o termo inter-grupo
\'{e} proporcional a \(\Phi > 0\); o termo intra-grupo \'{e}
proporcional a \(-\Omega\).

Esta estrutura justifica a forma de \(T_s^{(2),\text{sys}}\): o
numerador amplifica a sele\c{c}\~{a}o inter-grupo (coer\^{e}ncia,
associada a \(\Phi\langle A\rangle\)), o denominador amplifica a
sele\c{c}\~{a}o intra-grupo (extra\c{c}\~{a}o, associada a
\(\langle\Omega\rangle/\Phi\)).

\textit{Para a deriva\c{c}\~{a}o completa do marco evolutivo, veja §D.4.}



\section{\texorpdfstring{La Soglia di Sgretolamento \(T_s\)
(DEQ$^1$)}{La Soglia di Sgretolamento T\_s (DEQ$^1$)}}\label{la-soglia-di-sgretolamento-t_s-dequxb9}

La primeiro formulazione della DEQ, presentata nel \emph{Memoriale Tecnico}
(2026-04-23), pone:

\[
T_s^{(1)} = \frac{s_{iii} \cdot z \cdot \delta}{\sigma \cdot \varepsilon_{\text{dis}}}
\]

onde:

\begin{itemize}
\tightlist
\item
  \(s_{iii}\) è la \textbf{legitimidade ético-discursiva} del sistema
  (ereditata dalla Geometria);
\item
  \(z\) è la \textbf{proiezione sull'asse dialettico} (complessità del
  discorso pubblico --- eredità della Topografia);
\item
  \(\delta\) è l'\textbf{horizonte temporal} (Teorema Folk ---
  Geometria \S{}3);
\item
  \(\sigma\) è la \textbf{volatilidade estocástica} ereditata
  dall'estensione EDSD (Geometria \S{}12-13);
\item
  \(\varepsilon_{\text{dis}} \in [0,1]\) è il \textbf{coefficiente di
  disimpegno morale}, ispirato a Bandura (1999) e Milgram (1974).
\end{itemize}

\Proved{} \textbf{Lettura formale:} al numeratore stanno le forze di
coerenza (legitimidade $\times$ complessità $\times$ orizzonte); al denominatore le
forze di decoerência (volatilità $\times$ disimpegno). \(T_s^{(1)} = 1\) è la
limiar crítico. \(T_s^{(1)} \gg 1\) indica sistema stabile;
\(T_s^{(1)} < 1\) sistema in transição de fase imminente.

\Conj{} \textbf{Limite della DEQ$^1$:} l'equação \emph{non distingue} tra
sistemi semplicemente resilienti e sistemi antifragili, né tra sistemi
sotto attrito esogeno (volatilità esterna) e sistemi sotto attrito
endogeno (competizione involutiva). Due sistemi con lo stesso
\(T_s^{(1)}\) possono avere traiettorie radicalmente diverse. Questo
motiva l'estensione DEQ$^2$.



\section{I Due Assi Strutturali
Mancanti}\label{i-due-assi-strutturali-mancanti}

Introduciamo due variáveis che catturano dimensioni non incluse nella
DEQ$^1$.

\subsection{\texorpdfstring{L'Antifragilidade
\(A\)}{L'Antifragilidade A}}\label{lantifragilituxe0-a}

\Proved{} \textbf{Definizione formale (da Taleb 2012, formalizzazione
estrutural):}

\[
A = \left.\frac{\partial^2 \Pi}{\partial \sigma^2}\right|_{\sigma_0}
\]

onde \(\Pi\) è il payoff di sistema e \(\sigma_0\) è il livello base di
volatilità. \(A\) misura la \textbf{convessità della risposta allo
stress}:

\begin{itemize}
\tightlist
\item
  \(A > 0\): convexidade positiva $\Rightarrow$ il sistema \emph{guadagna} dalla
  volatilità (antifrágil);
\item
  \(A = 0\): linearità $\Rightarrow$ sistema resiliente neutro;
\item
  \(A < 0\): concavità $\Rightarrow$ sistema fragile (perde sproporzionalmente con
  lo stress).
\end{itemize}

\Hypoth{} \textbf{Proprietà:} \(A\) è una \emph{proprietà estrutural}
del sistema, non una risposta post-crisi. È determinata
dall'architettura del sistema (ridondanza, modularità, opzionalità)
primeiro che la crisi si manifesti.

\subsection{\texorpdfstring{Il Coefficiente di Involuzione
\(\Omega\)}{Il Coefficiente di Involuzione \textbackslash Omega}}\label{il-coefficiente-di-involução-omega}

\Proved{} \textbf{Definizione formale:}

\[
\Omega = \frac{\Delta E_{\text{dissipata}}}{\Delta E_{\text{utile}}} - 1
\]

onde \(\Delta E\) è il flusso di energia sistemica (risorse, capitale,
attenzione) per unità di tempo. Interpretazione:

\begin{itemize}
\tightlist
\item
  \(\Omega = 0\): ogni unità di energia dissipata produce un'unità di
  valore utile (break-even involutivo);
\item
  \(\Omega > 0\): il sistema dissipa più di quanto produca ---
  \textbf{Neijuan} in senso stretto;
\item
  \(\Omega \to \infty\): colapso energetico-produttivo.
\end{itemize}

\Hypoth{} \textbf{Ancoraggio empirico:} \(\Omega\) è mensurável come
complemento della Total Factor Productivity (TFP). Sistemi con TFP in
declino persistente hanno \(\Omega\) crescente.

\Hypoth{} \textbf{Parentela con Turchin:} la \emph{elite overproduction}
misurata come inverso del reddito relativo delle élite è un caso
particolare di \(\Omega\) applicato al sottosistema elitario. Turchin
cattura \(\Omega_{\text{élite}}\); la DEQ$^2$ generalizza.

\Proved{} \textbf{Microfondazione evolutiva: derivazione da CPGG.}
Quando la struttura di cattura è descrivibile por meio de il CPGG
(§D.4), \(\Omega\) è \emph{derivabile} dai parâmetros del gioco ---
non solo definito. All'equilibrio del CPGG, con frazione estrattiva
\(\phi\), quota asimmetrica \(\eta > 1/N\) e moltiplicatore \(r\) del
bene pubblico:

\[
\Omega(\phi,\eta,r) = \frac{\phi\cdot\eta\cdot r}{(1-\phi\eta)r - 1}
\]

Per \(\phi \ll 1\) (estrattori rari ma organizzati):

\[
\boxed{\Omega \approx \frac{\phi\cdot\eta\cdot r}{r-1}} \tag{P4}
\]

Il comportamento limite è rivelatore: per \(r \to 1^+\) (nessun
surplus cooperativo), \(\Omega \to \infty\) --- ogni estrazione è
totalmente distruttiva; per \(r \to \infty\), \(\Omega \to \phi\eta\)
(quota asimmetrica pura). \(\Omega\) non è una proprietà dell'attore
singolo ma una funzione del trio (\(\phi\), \(\eta\), \(r\)):
concentrazione estrattiva, asimmetria e produttività collettiva.



\section{\texorpdfstring{Derivazione di
\(T_s^{(2)}\)}{Derivazione di T\_s\^{}\{(2)\}}}\label{derivazione-di-t_s2}

\Hypoth{} \textbf{Principio di incorporazione:} \(A\) e \(\Omega\) non
sono variáveis indipendenti aggiuntive, ma \textbf{modulatori
strutturali} delle quantità già presenti in \(T_s^{(1)}\).

\textbf{Modulazione di \(\sigma\) por meio de \(A\).} Un sistema antifrágil
percepisce una volatilità effettiva ridotta, porque trasforma parte del
rumore in coerenza:

\[
\sigma_{\text{eff}} = \frac{\sigma}{1 + A}
\]

\begin{itemize}
\tightlist
\item
  \(A > 0\): \(\sigma_{\text{eff}} < \sigma\) (la volatilità è domata);
\item
  \(A = 0\): \(\sigma_{\text{eff}} = \sigma\) (caso neutro, recuperiamo
  la DEQ$^1$);
\item
  \(A \to -1^+\): \(\sigma_{\text{eff}} \to \infty\) (fragilidade
  catastrofica).
\end{itemize}

\textbf{Modulazione di \(\varepsilon_{\text{dis}}\) por meio de \(\Omega\).}
L'involução amplifica il disimpegno: ogni unità dissipata in
competizione improduttiva aumenta l'eteronomia del sistema
(Bandura/Milgram):

\[
\varepsilon_{\text{eff}} = \varepsilon_{\text{dis}} \cdot (1 + \Omega)
\]

\begin{itemize}
\tightlist
\item
  \(\Omega = 0\): nessuna amplificazione;
\item
  \(\Omega > 0\): disimpegno amplificato in proporzione all'involução.
\end{itemize}

\textbf{Equazione unificata DEQ$^2$ (forma individuale/meso):}

\[
\boxed{T_s^{(2)} = \frac{s_{iii} \cdot z \cdot \delta \cdot (1+A)}{\sigma \cdot \varepsilon_{\text{dis}} \cdot (1+\Omega)}}
\]

\Proved{} \textbf{Verifica di consistenza.} Per \(A = 0\) e
\(\Omega = 0\), \(T_s^{(2)} \to T_s^{(1)}\). La DEQ$^1$ è il caso limite
della DEQ$^2$ in assenza delle due modulazioni.



\section{Proprietà dell'Equazione}\label{proprietuxe0-dellequação}

\subsection{Monotonicità}\label{monotonicituxe0}

\Proved{} Per ogni variável, il segno della derivata è determinato:

\[
\frac{\partial T_s^{(2)}}{\partial s_{iii}},\ \frac{\partial T_s^{(2)}}{\partial z},\ \frac{\partial T_s^{(2)}}{\partial \delta},\ \frac{\partial T_s^{(2)}}{\partial A} > 0
\]

\[
\frac{\partial T_s^{(2)}}{\partial \sigma},\ \frac{\partial T_s^{(2)}}{\partial \varepsilon_{\text{dis}}},\ \frac{\partial T_s^{(2)}}{\partial \Omega} < 0
\]

Le prime sono \textbf{forze di coerenza}, le seconde \textbf{forze di
decoerência}. Nessuna variável ha segno ambiguo --- proprietà importante
per l'interpretabilità.

\subsection{Asintoti Critici}\label{asintoti-critici}

\begin{itemize}
\tightlist
\item
  \(\delta \to 0\): colapso dell'horizonte temporal (miopia totale) $\Rightarrow$
  \(T_s^{(2)} \to 0\).
\item
  \(\sigma \to \infty\): volatilità illimitata $\Rightarrow$ \(T_s^{(2)} \to 0\), a
  meno che \(A \to \infty\) (caso teorico di antifragilidade infinita).
\item
  \(\Omega \to \infty\): involução totale $\Rightarrow$ \(T_s^{(2)} \to 0\).
\item
  \(A \to -1\): fragilidade catastrofica $\Rightarrow$ \(T_s^{(2)} \to 0\) anche con
  volatilità finita.
\end{itemize}

\subsection{Punto Critico di Transizione di
Fase}\label{punto-critico-di-transizione-di-fase}

\Hypoth{} \textbf{Ipotesi operativa:} la transição de fase del sistema
(colapso, rivoluzione, riorganizzazione estrutural) si manifesta
quando \(T_s^{(2)} < 1\) per una durata persistente
\(\Delta t > \Delta t^*\), onde \(\Delta t^*\) dipende dalla scala del
sistema:

{\def\LTcaptype{none} % do not increment counter
{\small\begin{longtable}[]{@{}ll@{}}
\toprule\noalign{}
Scala & \(\Delta t^*\) stimato \\
\midrule\noalign{}
\endhead
\bottomrule\noalign{}
\endlastfoot
Individuo / azienda & 4-12 mesi \\
Sistema nazionale & 3-7 anni \\
Ciclo egemonico & 15-25 anni (coerente con Arrighi) \\
\end{longtable}}
}

\Conj{} \textbf{Congettura di scala:} \(\Delta t^*\) è
approssimativamente proporzionale alla radice del numero di attori del
sistema --- hipótese da verificaçãore empiricamente nei capitoli applicati
(Parte III).



\section{Limiti di Questa
Formulazione}\label{limiti-di-questa-formulazione}

\Conj{} \textbf{Limite 1.} \(T_s^{(2)}\) è qui scritto a \emph{livello
local/meso} --- descrive un attore come se le sue componenti fossero
allineate (\(\Phi = 1\)). L'aggregazione formale al sistema globale, già
anticipata nel \S{}3.1bis por meio de l'operatore \(\langle\cdot\rangle_\Phi\),
è completata nel Cap.~4 onde \(T_s^{(2),\text{sys}}\) assume la sua
forma definitiva con \(\Phi < 1\). La differenza tra \(T_s^{(2)}\) e
\(T_s^{(2),\text{sys}}\) misura quanto un sistema \emph{spreca}
legitimidade per decoerência interna --- il caso paradigmatico è l'USA
contemporaneo.

\Conj{} \textbf{Limite 2.} Le variáveis
\(s_{iii}, z, \sigma, \varepsilon_{\text{dis}}, A, \Omega\) sono qui
assunte misurabili. La definição operativa di \emph{come misurarle} è
rinviata all'\textbf{Apêndice B} (Protocolli di medição dei
parâmetros).

\ToVerify{} \textbf{Da verificaçãore:} la scelta della forma moltiplicativa
\((1+A)\) e \((1+\Omega)\) è semplice e monotonicamente consistente, ma
non è l'unica possibile. Forme alternative (esponenziale, razionale)
andrebbero testate su dati storici. Questa verificação è rinviata al paper
tecnico in inglese.



\section{Síntese do Capítulo}\label{sintesi-del-capítulo}

\Proved{} \textbf{In una frase:} la stabilità di un sistema egemonico è
il rapporto tra una coerenza (legitimidade $\times$ dialettica $\times$ orizzonte)
\emph{amplificata} dall'antifragilidade e una decoerência (volatilità $\times$
disimpegno) \emph{amplificata} dall'involução.

\[
T_s^{(2)} = \frac{\underbrace{s_{iii} \cdot z \cdot \delta}_{\text{coerenza}} \cdot \overbrace{(1+A)}^{\text{amplificatore}}}{\underbrace{\sigma \cdot \varepsilon_{\text{dis}}}_{\text{decoerenza}} \cdot \underbrace{(1+\Omega)}_{\text{amplificatore}}}
\]

Il Cap.~4 affronterà il passaggio al sistema aggregato introducendo il
parâmetro de ordem \(\Phi\).



\section{Notas Epistêmicas de Encerramento
Capítulo}\label{note-epistemiche-di-chiusura-capítulo}

{\def\LTcaptype{none} % do not increment counter
{\small\begin{longtable}[]{@{}
  >{\raggedright\arraybackslash}p{(\linewidth - 4\tabcolsep) * \real{0.3333}}
  >{\raggedright\arraybackslash}p{(\linewidth - 4\tabcolsep) * \real{0.3333}}
  >{\raggedright\arraybackslash}p{(\linewidth - 4\tabcolsep) * \real{0.3333}}@{}}
\toprule\noalign{}
\begin{minipage}[b]{\linewidth}\raggedright
\#
\end{minipage} & \begin{minipage}[b]{\linewidth}\raggedright
Afirmação
\end{minipage} & \begin{minipage}[b]{\linewidth}\raggedright
Status
\end{minipage} \\
\midrule\noalign{}
\endhead
\bottomrule\noalign{}
\endlastfoot
1 & La correzione notazionale di \(\delta\) & \Proved{} coerenza
interna \\
2 & L'identità formale tra \(\Pi\) Topografia e \(\Pi\) Geometria
(\(\delta^*=0{,}50\)) & \Proved{} verificaçãodo sui due testi \\
3 & L'operatore di trasposizione \(\langle\cdot\rangle_\Phi\) (Kuramoto)
& \Hypoth{} hipótese formale di scala \\
4 & La derivazione di \(T_s^{(1)}\) dal Memoriale & \Proved{} richiamo
fedele \\
5 & Le definizioni di \(A\) e \(\Omega\) & \Proved{} formalizzazioni
rigorose \\
6 & L'equação \(T_s^{(2)}\) & \Hypoth{} hipótese operativa
unificante \\
7 & La condizione \(T_s^{(2)} < 1\) come soglia di transizione &
\Hypoth{} hipótese testabile \\
8 & La scala \(\Delta t^*\) & \Conj{} conjetura di scala \\
9 & La forma moltiplicativa \((1+A)\), \((1+\Omega)\) & \ToVerify{} da
verificaçãore empiricamente \\
\end{longtable}}
}


